% Track d7 -- The b-field between manifolds.
% Verification: investigations/rest_energy/verify/d7_bfield_between.py (claims d7-1a ... d7-8).
% New references are in sections/references.tex.
\section{\texorpdfstring{The $b$-field between manifolds}{The b-field between manifolds}}
\label{d7:sec:main}

The author's first remark is that the $b$-particle is a deal breaker when two manifolds or more are in question. \Cref{d5:thm:bfield2}
tested it for two binding frames and one field, the copy rotation. This section asks the question for any number of manifolds and any
field that can act between them, and answers it in five steps. First, which fields exist depends on how the Lorentz algebra acts: a common
action admits the complex orthogonal copy algebra $\so{n,\C}$, a spectator action decouples the Lorentz-carrying manifold from every
binding-compatible field (fixed, partly restated from \cite{ES17}; \cref{d7:sec:exist}). Second, the closure itself admits such a field: with a
common Lorentz action the rotor may act by any element of $\uu n$, which is a copy rotation plus a matrix of rates (constructed;
\cref{d7:sec:rotor}). Third, every binding-compatible field turns the frame flow into a Schr\"odinger flow on $\C^{3n}$, unitary for the
undivided energy exactly when the field is compact. The rates are gauge invariants, the field redistributes the frequencies at fixed sum, and
it is removable exactly when it preserves the rate eigenspaces (fixed; \cref{d7:sec:schr,d7:sec:boost}). Fourth, three manifolds carry an
invariant that two do not: the sign of $B_{12}B_{23}B_{31}$, odd under interior time reversal, the analogue of Jarlskog's invariant (fixed;
\cref{d7:sec:three}). Fifth, the copy isospin is integral: frames carry isospin one, never one-half. The Yang--Mills doublet acting on two copies
of the undivided energy is a connection in one of its three directions and a rate term in the other two, and doublets exist only under the
spectator action, where unequal internal rates break $\so4_\Delta$ of \cite{ES17} to a diagonal $\so3$ (fixed; \cref{d7:sec:isospin}). The verdict
is in \cref{d7:sec:assess}. Identifiers d7-$n$ refer to the \texttt{PASS} lines of \texttt{verify/d7\_bfield\_between.py} (\cref{d7:tab:claims}).

\paragraph{Notation.} As in \cref{d5:sec:main}, with $n$ copies: $V_n=\R^n\otimes\R^6$ in the copy-major basis, $\eta_n=\id_n\otimes\eta$,
$G_n=\id_n\otimes G$, and $z=(z_1,\dots,z_n)\in\C^{3n}$ with $z_k=\bB_k+i\bE_k$, on which $G_n$ acts as $i$. $E_{kl}$ are the matrix units on the copy
index. The \emph{copy rotation} $C^{kl}=(E_{lk}-E_{kl})\otimes\id_6$ maps copy $k$ to copy $l$, and the \emph{copy boost} is $(E_{lk}-E_{kl})\otimes G$. A
\emph{copy field} is $b=B\otimes\id_6$ with $B\in\so n$ real antisymmetric, and $J:=B_{12}B_{23}B_{31}$ at $n=3$. The binding term with
rates $\kappa_k\ge0$ is $K\otimes G$ with $K=\operatorname{diag}(\kappa_1,\dots,\kappa_n)$, and $S=K\otimes G+b$, $\dot X=-SX+\pi/m$, $\dot\pi=S^\dagger\pi$, with
$M$ as in \cref{d5:thm:bfield2}. A field is \emph{binding-compatible} if it commutes with $G_n$. Rates are rates, and reading one as a rest
energy needs $mc^2=\hbar\kappa$ (\cref{d3:prop:norel}).

\subsection{Which fields can act between manifolds}
\label{d7:sec:exist}

\begin{proposition}[Inter-manifold fields and the Lorentz action]
\label{d7:prop:exist}
\begin{enumerate}[label=(\roman*),leftmargin=2.2em]
\item \emph{Common binding action} (the diagonal $\lG$). The elements of $\so{3n,3n}$ that commute with it and with $G_n$ are
$\{B\otimes\id_6+B'\otimes G:\ B,B'\in\so n\}\cong\so{n,\C}$, of dimension $0,2,6,12$ for $n=1,\dots,4$. The compact part consists of the copy
rotations, of dimension $0,1,3,6$; the copy boosts $B'\otimes G$ are noncompact.
\item \emph{Common interior action} (the diagonal $\lint$, the action $\Delta$ of Section~10 of \cite{ES17}). The elements of
$\so{3n,3n}$ commuting with it and with $G_n$ are the copy rotations alone, of dimension $0,1,3$ for $n=1,2,3$; no copy boost commutes with $\lint$.
At $n=2$ the copy rotation is $C^{12}=M_{12}+C_{\rm surp}$, where $M_{12}=\varepsilon\otimes P_{\rm int}$ is the copy rotation of the interior
four-planes, the copy algebra of Theorem~10.2 of \cite{ES17}, and $C_{\rm surp}=\varepsilon\otimes P_{\rm surp}$ that of the surplus pairs. Both
commute with the diagonal $\lint$, and neither commutes with $G_2$.
\item \emph{Spectator action} ($\lG^{(1)}$ or $\lint^{(1)}$ on copy~$1$ alone). The elements of $\so{3n,3n}$ commuting with it and with $G_n$ form
$\so{3n-3,\C}$ of copies $2,\dots,n$, of dimension $6$ and $30$ at $n=2,3$. Every one of them vanishes on copy~$1$ and has no component into it.
\end{enumerate}
\end{proposition}

Part (i) restates the third row of Proposition~13.14 of \cite{ES17} (fixed there) and \cref{d5:prop:so6C}(iii) at $n=2$. Part (iii) restates
Theorem~17.20 of \cite{ES17} at $n=3$ (fixed there), and \cref{d5:prop:so6C}(iv) at $n=2$. Its proof is Schur's lemma: $V$ is an irreducible,
nontrivial module of the Lorentz algebra of copy~$1$, and the other copies are trivial modules of it. Part (ii) and the decomposition of
$C^{12}$ are fixed (d7-1a to d7-1d). Two consequences matter below. A binding-compatible, Lorentz-invariant field can couple the manifold that
carries the Lorentz algebra to the others only when that algebra acts on all of them in the same way. Under the spectator action copy~$1$ is out
of reach, and the inter-manifold fields act among the internal manifolds only. And the paper's copy algebra $\mathfrak k_{\rm cop}$ does not respect
the binding: the binding-compatible copy rotation is $M_{kl}$ plus the surplus copy rotation.

\subsection{The closure admits an inter-manifold field}
\label{d7:sec:rotor}

Theorem~17.7 of \cite{ES17} shows that, on one constituent, the only equivariant closure is $\R\ltimes_GV$: the rotor acts by $G$. After it the
paper extends the result to several constituents ``factor by factor, since equivariant maps between distinct constituents vanish'', with a common
rate. That is correct for a Lorentz algebra acting on each constituent separately. For a common action it is not the only possibility.

\begin{proposition}[Rotor generators on $n$ binding frames]
\label{d7:prop:rotor}
The endomorphisms of $V_n$ commuting with the diagonal $\lG$ form $\mathfrak{gl}(n,\C)=\mathfrak{gl}(n)\otimes\operatorname{span}\{\id,G\}$, of dimension
$2n^2$. Those that rotate the frame covector without dilating it, the condition of Theorem~17.7 of \cite{ES17}, form
\[
\uu n=\so n\otimes\id_6\ \oplus\ \operatorname{Sym}(n)\otimes G ,
\]
of dimension $n^2$. The first summand is $\eta_n$-antisymmetric (copy rotations) and the second is $\eta_n$-self-adjoint (rate matrices). The paper's
rotor $G_n$ spans the centre. For $\mathfrak s_n=\R\ltimes_AV_n$ with $A=K\otimes G+B\otimes\id_6\in\uu n$ and the rotor clamped at $\theta=\Omega\tau$, the
substitution $X=e^{-\theta A}x$ of Section~17.3 of \cite{ES17} gives $e^{-\theta A}\dot x=\dot X+\Omega AX$, so the shift is $S=\Omega A$.
\end{proposition}

At $n=2$ the algebra is the energy symmetry $\uu2$ of \cref{d5:prop:energysym}(i). The dimensions and the decomposition are fixed (d7-2a); the
reading as rotor generators is new. The semidirect product is a Lie algebra for every $A$, so the closure exists. That a
non-central $A$ is admissible is constructed: nothing in the algebra selects it, and Theorem~17.7 of \cite{ES17} selects $G_n$ only once the
action is taken per constituent, or a common rate is imposed. With a non-central $A$ the copy field and unequal rates enter the closure
together, as the two halves of one generator. Their ratios are fixed by $A$ and their scale by the rotor. This is the one place found where the
arena itself supplies the unequal rates that \cref{d5:thm:bfield2} needs. The paper's one-rotor closure, with $A=G_n$, gives every constituent the
same rate (Proposition~17.18 of \cite{ES17}).

\subsection{The frame flow is a Schr\"odinger flow}
\label{d7:sec:schr}

\begin{theorem}[Binding-compatible fields between $n$ manifolds]
\label{d7:thm:schr}
Let $b\in\Cent(G_n)\cap\so{3n,3n}$. By \cref{d5:prop:so6C}(i), $b$ acts on $z$ as $R+iQ$ with $R,Q\in\so{3n}$ real: $R$ is
its paired-rotation part and $Q$ its paired-boost part. Put $\tilde K=K\otimes\id_3$ and $H:=\tilde K-iR+Q$.
\begin{enumerate}[label=(\roman*),leftmargin=2.2em]
\item $S$ acts on $z$ as $iH$, so the free frame flow is $z(\tau)=e^{-iH\tau}z(0)$.
\item The flow preserves $E_6=\tfrac12|z|^2$ for every $K$ and every initial value iff $Q=0$, that is, iff $b$ is compact; then $H$ is Hermitian. The
compact fields, the only ones admissible as Yang--Mills fields (Corollary~2.4 of \cite{ES17}), are exactly those whose frame flow is unitary for the
undivided energy.
\item The rates are gauge invariants. Under any $\eta_n$-isometric change of frame $X=U(\tau)Y$, the $\eta_n$-self-adjoint part of $S$ is conjugated,
so its spectrum $\{\pm i\kappa_k\}$ is unchanged. No field, compact or not, changes a rate.
\item For compact $b$: $\tr H=3\sum_k\kappa_k$ and $\tr H^2=3\sum_k\kappa_k^2+\|R\|^2$. The field redistributes the frequencies of the frame at fixed
sum and widens their spread (level repulsion).
\item $b$ is removable along the worldline iff $[\tilde K,R+iQ]=0$, that is, iff it preserves every rate eigenspace. At $n=3$ the compact
binding-compatible fields, of dimension $36$, have a removable part of dimension $36$, $18$ and $9$ for rates $(\kappa,\kappa,\kappa)$, $(\kappa,\kappa,\kappa')$
and three distinct rates. Within the copy $\so3$ the dimensions are $3$, $1$ and $0$.
\item (Boundedness) Let $b$ be compact. Every orbit of $(X,\pi)$ is bounded iff, for every $\lambda$ with $\lambda,-\lambda\in\spec H$ (including $\lambda=0$),
the eigenspaces $\mathcal E_\lambda(H)$ and $\mathcal E_{-\lambda}(H)$ are orthogonal for the bilinear form $v\cdot w=\sum_av_aw_a$ on $\C^{3n}$.
\end{enumerate}
\end{theorem}

\begin{proof}
(i) On the sector $G_n=i$, $K\otimes G$ acts as $i\tilde K$, a paired rotation as $R$ and a paired boost as $iQ$ (\cref{d4:prop:lp,d5:prop:so6C}); $\dot X=-SX$ becomes $\dot z=-iHz$. (ii) $E_6=\tfrac12|X|^2$, so it is conserved for all $X$ iff $S$ is Euclidean
antisymmetric. $K\otimes G$ and the paired rotations are, and the paired boosts are Euclidean symmetric. (iii) is
\cref{d4:thm:covariant}(ii) on $V_n$. The self-adjoint part transforms as a tensor, $U^{-1}(K\otimes G)U$, while $U^{-1}bU+U^{-1}\dot U$ stays in
$\so{3n,3n}$. (iv) $\tr(\tilde KR)=0$ because $\tilde K$ is symmetric and $R$ antisymmetric, and $\tr(-iR)^2=\|R\|^2$. (v) is
\cref{d4:prop:larmor}: $e^{-b\tau}$ intertwines the systems iff $b$ commutes with $K\otimes G$. (vi) is \cref{d4:thm:criterion} with
$S^{\!\top}=-S$ (both summands are Euclidean antisymmetric). Then $\mathcal E_\mu(S^\dagger)=\eta_n\mathcal E_{-\mu}(S)$, and the condition reads
$\eta_n(\mathcal E_\nu(S),\mathcal E_{-\nu}(S))=0$ for every $\nu$. The sectors $G_n=\pm i$ are $\eta_n$-orthogonal, and on the sector $G_n=i$,
$\eta_n=2\,v\cdot w$ in the coordinates $z$. The eigenspaces of $S$ there are those of $iH$, and on the other sector their conjugates.
\end{proof}

Parts (i), (ii), (iv) and (v) are fixed but elementary: the arena's version of the fact that a compact gauge field acts unitarily
(d7-3a, d7-3b, d7-3d, d7-3e). Part (iii) is \cref{d4:thm:covariant}(ii) restated for $n$ copies, checked exactly on a Cayley curve (d7-3c). Part (vi) is
a corollary of \cref{d4:thm:criterion}. On five cases at $n=3$, exact semisimplicity of $M$, the criterion and the growth of $\|e^{Mt}\|$ agree
(d7-3f). Bounded: rates $(1,2,4)$ with $(B_{12},B_{23},B_{31})=(\tfrac12,\tfrac13,\tfrac15)$; the seesaw $(0,3,5)$ with $(1,\tfrac12,\tfrac13)$, a
frame without rate confined by two couplings; equal rates $(1,1,1)$ with $B_{12}=1$. Unbounded: $(1,2,3)$ with $(1,1,1)$, where $\spec H=\{0,2,4\}$ and the
zero mode has $v\cdot v\neq0$, a resonance created by the field between manifolds with different rates; and zero rates with $(1,1,1)$, where
the field produces only drift. Part (vi) generalises \cref{d5:thm:bfield2}(v): at $n=2$, $\lambda_++\lambda_-=\kappa_1+\kappa_2$, so the only coincidence
is $\lambda_-=0$, at $\beta^2=\kappa_1\kappa_2$.

\subsection{Copy boosts: metric connections, not gauge fields}
\label{d7:sec:boost}

\begin{proposition}[The copy boost between two manifolds]
\label{d7:prop:boost}
Let $n=2$ and $b=\beta'\varepsilon\otimes G$, so that $H=K+\beta'\varepsilon$ on the copy index (\cref{d7:thm:schr}(i)), real and not symmetric. Its eigenvalues are
$\tfrac12(\kappa_1+\kappa_2)\pm\sqrt{\tfrac14(\kappa_1-\kappa_2)^2-\beta'^2}$.
\begin{enumerate}[label=(\roman*),leftmargin=2.2em]
\item Below the exceptional point, $|\beta'|<\tfrac12|\kappa_1-\kappa_2|$, the spectrum is real, $M$ is semisimple with imaginary spectrum, and every orbit
is bounded. The flow then preserves a positive form $z^\dagger\mathcal Qz$ with $H^\dagger\mathcal Q=\mathcal QH$, which is not $E_6$.
\item At $|\beta'|=\tfrac12|\kappa_1-\kappa_2|$, $M$ is not semisimple; above it, $M$ has eigenvalues with real part $\pm\sqrt{\beta'^2-\tfrac14(\kappa_1-\kappa_2)^2}$.
\item At equal rates the threshold is zero: every $\beta'\neq0$ grows exponentially, although $b$ commutes with $K\otimes G$ and is removable, by a
boosted frame $e^{-b\tau}$ that is itself unbounded.
\end{enumerate}
\end{proposition}

This is fixed, checked exactly at rates $(1,4)$ with $\beta'=\tfrac12,1,\tfrac32,2$, and at $(2,2)$ with $\beta'=\tfrac13$ (d7-4a, d7-4b). It is the
copy analogue of the pairing boosts of Proposition~16.7 of \cite{ES17} and of \cref{d4:sec:pG}. The spectral pattern is inherited: the
non-Hermitian dimer with an exceptional point \cite{Kato66,BenderBoettcher98}, and the conserved metric of a pseudo-Hermitian Hamiltonian
\cite{Mostafazadeh02}. A copy boost is a metric connection, legitimate in the interior Lagrangian, but not a gauge field (Corollary~2.4 of
\cite{ES17}). Where it keeps the motion bounded it does so at the price of the undivided energy, which it replaces by another positive form.

\subsection{Three manifolds: a time-odd invariant}
\label{d7:sec:three}

\begin{theorem}[The copy field between three manifolds]
\label{d7:thm:three}
Let $n=3$, distinct rates, and $b=B\otimes\id_6$, so $H=(K-iB)\otimes\id_3$.
\begin{enumerate}[label=(\roman*),leftmargin=2.2em]
\item $\det(x-K+iB)=x^3-(\textstyle\sum\kappa)x^2+(\sum_{k<l}\kappa_k\kappa_l-\sum B^2)x-(\kappa_1\kappa_2\kappa_3-\kappa_3B_{12}^2-\kappa_1B_{23}^2-\kappa_2B_{31}^2)$: the spectrum
depends on the couplings only through their squares.
\item The plaquette $H_{12}H_{23}H_{31}=iJ$ is purely imaginary, and $\det[H,K]=2iJ(\kappa_1-\kappa_2)(\kappa_1-\kappa_3)(\kappa_2-\kappa_3)$.
\item The endomorphisms commuting with $G_3$, the diagonal $\lG$ and $K\otimes G$ are $\sum_k(a_k+b_kG)$ on copy $k$, and the $\eta_3$-isometries among them are
exactly the eight copy reflections $\operatorname{diag}(\pm1,\pm1,\pm1)\otimes\id_6$. Under them $B_{kl}\mapsto s_ks_lB_{kl}$, so the invariants of the field
are $|B_{12}|,|B_{23}|,|B_{31}|$ and the sign of $J$.
\item With $P=\id_3\otimes\operatorname{diag}(\id_3,-\id_3)$ ($\bE\mapsto-\bE$) and $T=-P$ ($\bB\mapsto-\bB$), parity maps $(K,B)$ to $(-K,B)$, and interior
time reversal, $X(\tau)\mapsto TX(-\tau)$, maps $(K,B)$ to $(K,-B)$. At $n=2$ a copy reflection undoes $B\mapsto-B$; at $n=3$ with $J\neq0$ none does, since
$J\mapsto-J$. For general $n$, $(K,B)$ and $(K,-B)$ are related by a copy reflection iff the graph of the nonzero couplings is bipartite.
\item With energy started in copy~$1$, the copy energies $P_k(\tau)=|z_k(\tau)|^2/|z(0)|^2$ obey $P_k(\tau;-B)=P_k(-\tau;B)$. $P_1$ is the same for $B$ and
$-B$, and $P_2,P_3$ are not. At rates $(1,2,4)$ with $B_{12}=B_{23}=B_{31}=1$ and $\tau=\tfrac12$, $P_2=0.293$ for $B$ against $0.116$ for $-B$: the sign
of $J$ is the sense in which the energy circulates around the three manifolds.
\end{enumerate}
\end{theorem}

\begin{proof}
(i), (ii) Symbolic expansion (d7-5a, d7-5b). For a $3\times3$ Hermitian matrix, $\det H$ contains $2\Re(H_{12}H_{23}H_{31})$, which vanishes here.
(iii) The commutant is exact (d7-5c). $(a+bG)^{\!\top}\eta(a+bG)=(a^2-b^2)\eta+2ab\,\eta G$ because $G^{\!\top}\eta=\eta G$ and $G^2=-\id$, so $ab=0$
and $a^2-b^2=1$. (iv) $PGP^{-1}=-G$, $P$ commutes with $B\otimes\id_6$, and $\dot X=-SX$ goes to $\dot X'=-S'X'$ with $S'=-TST^{-1}$. A copy
reflection $s$ maps $B$ to $-B$ iff $s_ks_l=-1$ on every edge, which is a two-colouring (d7-5d, d7-5e). (v)
$|(e^{-i\bar H\tau})_{k1}|=|(e^{iH\tau})_{k1}|$, and $P_1$ is even in $\tau$ for Hermitian $H$ (d7-5f).
\end{proof}

The statements are fixed. Their pattern is inherited three times over. With two families there is no rephasing-invariant phase and with three
there is one, the Kobayashi--Maskawa phase \cite{KobayashiMaskawa73}, whose invariant measure is Jarlskog's commutator determinant
\cite{Jarlskog85}, the form of (ii). A purely imaginary hopping around a triangle is a flux $\pm\pi/2$ through the plaquette, the Peierls phase
\cite{Peierls33} of a lattice gauge field. And time-reversal-odd hopping makes a quantum walk chiral \cite{Zimboras13}. What the arena adds is
the identification. The copy field is an $\eta$-compatible gauge field on the three-site ``lattice'' of manifolds, with $J$ its Wilson loop. Its
phase is fixed at $\pm\pi/2$ because the field is real: per-copy duality rotations, which would move it, are not $\eta$-isometries
(Proposition~\ref{d3:prop:norel}(ii)). On one manifold the binding is odd under parity and even under interior time reversal. At two manifolds the
copy field is even under both, up to relabelling. From three on, a field whose couplings close a triangle is odd under interior time reversal,
and no relabelling removes it.

\subsection{Isospin one, not one-half}
\label{d7:sec:isospin}

Yang and Mills' $b$-field acts on the isotopic doublet of nucleons \cite{YangMills54}. The question is whether an inter-manifold field can.

\begin{proposition}[Copy isospin]
\label{d7:prop:isospin}
\begin{enumerate}[label=(\roman*),leftmargin=2.2em]
\item On the tower module, the spin lift of the copy rotation of a set $D$ of directions is $-\tfrac12\sum_{a\in D}\eta_{aa}e^{(k)}_ae^{(l)}_a$, a sum of $|D|$
commuting terms of square $-\tfrac14$, so $\exp(2\pi\,\cdot)=(-1)^{|D|}$. At $n=2$: for the six binding directions the eigenvalues are $i\{-3,\dots,3\}$ and
$\exp(2\pi C)=+\id$. For the four interior directions ($M_{12}$) they are $i\{-2,\dots,2\}$, $+\id$. For the three temporal directions they are
$i\{\pm\tfrac12,\pm\tfrac32\}$, $-\id$, and that rotation does not commute with $G_2$. A binding-compatible copy rotation pairs each temporal direction with
its partner, so $|D|$ is even.
\item At $n=3$ the three binding copy rotations close on $\so3$ on the $512$-dimensional module, with Casimir $j(j+1)=0,2,6,12$ on $40,216,200,56$
dimensions ($j=0,\dots,3$), and $\exp(2\pi C^{12})=+\id$. The frames form the vector, $j=1$.
\item The doublet $\uu2$ acting on $(z_1,z_2)$ preserves $E_6\oplus E_6$ and commutes with $G_2$ and the diagonal $\lG$. In real form $i\id\mapsto\id\otimes G$,
$i\sigma_x\mapsto\sigma_x\otimes G$, $i\sigma_y\mapsto-\varepsilon\otimes\id_6$ and $i\sigma_z\mapsto\sigma_z\otimes G$. Only $i\sigma_y$, the copy rotation, is
$\eta_2$-antisymmetric. The other three are $\eta_2$-self-adjoint, terms of the same type as the binding. For $S=\kappa G_2+b_{\rm YM}$,
$b_{\rm YM}=\sum_a\beta_a\,(i\sigma_a)$, the lines are $\kappa\pm|\boldsymbol\beta|$, and the rates, the spectrum of the self-adjoint part, are
$\kappa\pm(\beta_x^2+\beta_z^2)^{1/2}$.
\end{enumerate}
\end{proposition}

This is fixed (d7-6a to d7-6c). Part (i) restates Proposition~10.9 of \cite{ES17} (``copy rotations act through orthogonal groups'') and
\cref{d5:prop:spin}(i), and adds that the binding is what makes the lift single-valued. Under a common Lorentz action the inter-manifold field
is therefore an $SO(n)$ gauge field. At $n=3$ it acts on triplets, as in the model of Georgi and Glashow \cite{GeorgiGlashow72}, and never on
doublets. Part (iii) restates \cref{d5:prop:energysym}(i) in the Yang--Mills reading and adds the spectra; it is the precise form of
``Yang--Mills between two copies''. The doublet group exists as a symmetry of the undivided energy
(the copy $\uu2$), but $\eta_2$ cuts it to its circle. The two directions it removes are not connections at all:
they belong to the symmetric part, where the binding $c^2$ lives in the dictionary of \cref{s0:sec:question}. A doublet field that splits the
rates between two manifolds does so only through those two directions. In the arena that is a modification of the symmetric part, not a
$b$-field.

\begin{proposition}[The weak isospin of the paper between unequal constituents]
\label{d7:prop:so4unequal}
At $n=3$, with the numbering of Section~17 of \cite{ES17} and $T^\pm_i$ as in \cref{d4:prop:PS}, let $K=\operatorname{diag}(\kappa_1,\kappa_2,\kappa_3)$ with
$\kappa_2\neq\kappa_3$.
\begin{enumerate}[label=(\roman*),leftmargin=2.2em]
\item The part of $\so4_\Delta$ commuting with $K\otimes G$ is $\operatorname{span}\{\mathrm{rot}(13,14),\mathrm{rot}(13,15),\mathrm{rot}(14,15)\}
=\operatorname{span}\{T^+_1+T^-_1,\ T^+_2-T^-_2,\ T^+_3-T^-_3\}$, the graph of an isomorphism $\su2_+\to\su2_-$: a diagonal $\so3$. The three directions
$\mathrm{rot}(7,13),\mathrm{rot}(7,14),\mathrm{rot}(7,15)$, which join constituent two to constituent three, are not removable.
\item For $\kappa=(5,3,5)$ and $b=2T^+_3$,
\[
\det(\lambda-S)=(\lambda^2+9)^2(\lambda^2+16)(\lambda^2+25)^3(\lambda^2+36)(\lambda^4+36\lambda^2+196),
\]
and $M$ is semisimple. The quartic carries the lines $4\pm\sqrt2$, the two-level mixing of $z_7$ and $z_{13}$ with rates $3,5$ and coupling $1$
(\cref{d5:thm:bfield2}(ii)). At equal rates $5$ the same field gives $(\lambda^2+25)^5(\lambda^2+16)^2(\lambda^2+36)^2$, the removable Zeeman pattern
of \cref{d4:prop:so4spec}(i).
\end{enumerate}
\end{proposition}

This is fixed (d7-7a, d7-7b). It completes \cref{d4:prop:so4spec}, which used equal rates. The paper's binding-compatible non-abelian algebra is
itself an inter-manifold field: four of its six generators join constituent two to constituent three. It carries doublets ($(\tfrac12,\tfrac12)$ on
$W$) because the spectator action leaves constituents two and three free of the Lorentz algebra (\cref{d7:prop:exist}(iii)). With equal rates it is
removable along a worldline, as \cref{d4:sec:three} found. With unequal rates it breaks to a diagonal $\so3$, and its three remaining directions
become physical. The pattern $\su2\oplus\su2\to\su2_{\rm diag}$ is that of chiral symmetry breaking and of the custodial symmetry of the Higgs
sector \cite{SikivieEtAl80}. That is an interpretive parallel only: here the breaking is by a constant rate matrix, not by a vacuum.

\subsection{Assessment}
\label{d7:sec:assess}

\begin{remark}[The deal breaker, settled]
\label{d7:rem:verdict}
\begin{enumerate}[label=(\alph*),leftmargin=2.2em]
\item \emph{True, in this sense} (fixed). Between manifolds with unequal rates, a compact binding-compatible field cannot be removed. It shifts and
splits the lines, it can confine a frame that has no rate of its own, it can create a resonance, and from three manifolds on it is non-abelian and
carries a time-odd invariant. The one-manifold no-go of \cref{d4:sec:bfield} does not extend to $n\ge2$
(\cref{d5:thm:bfield2,d7:thm:schr,d7:thm:three}).
\item \emph{False, in this sense} (fixed). The field never creates or changes a rate. Rates are the spectrum of the self-adjoint part, a tensor under
every change of frame. The field only redistributes the frequencies, at fixed sum. With equal rates every binding-compatible field, abelian or
the paper's $\so4_\Delta$, is removable along a worldline. With zero rates it produces only drift (\cref{d7:thm:schr}(iii)--(vi)).
\item \emph{Where unequal rates come from.} Not from the paper's closure: its single rotor acts through $G_n$ at one rate, so within \cite{ES17} as it stands
the inter-manifold field is removable (fixed for that closure). Under a common Lorentz action the closure admits a rotor acting by a non-central element
of $\uu n$. The copy field and the rate differences then arrive together, as two halves of one generator (constructed: admitted, not
selected; \cref{d7:prop:rotor}).
\item \emph{Isospin one-half} (fixed). Under a common Lorentz action the copy isospin is integral, and the doublet field on two copies is a connection
in one of its three directions only. Doublets need the spectator action, where the Lorentz-carrying manifold is decoupled
(\cref{d7:prop:exist,d7:prop:isospin,d7:prop:so4unequal}).
\item \emph{The arrow reversed} (interpretive). The non-removable part of the copy field is $[K\otimes G,B\otimes\id_6]=[K,B]\otimes G$, with
$[K,B]_{kl}=(\kappa_k-\kappa_l)B_{kl}$ and $\|[K\otimes G,b]\|^2=6\sum_{k\neq l}(\kappa_k-\kappa_l)^2B_{kl}^2$ (d7-8, fixed). Formally this is the gauge-boson
mass term $g^2\|[T,\langle\Phi\rangle]\|^2$ of the Englert--Brout--Higgs--Kibble mechanism \cite{EnglertBrout64,Higgs64,Kibble67}, with the rate matrix
in the role of an adjoint vacuum value. The arena thus reverses the arrow of the hypothesis. It is not the $b$-field that differentiates the rates;
the differentiation of the rates, which breaks the copy symmetry, is what gives the $b$-field an effect. The parallel is not dynamical: here $b$ is a
constant connection without kinetic term, and the rates are inputs.
\end{enumerate}
\end{remark}

\begin{table}[t]
\centering
\footnotesize
\begin{tabular}{@{}>{\raggedright\arraybackslash}p{8.1cm}>{\raggedright\arraybackslash}p{5.0cm}>{\raggedright\arraybackslash}p{2.0cm}@{}}
\toprule
statement & class & checks\\
\midrule
\cref{d7:prop:exist}(i), (iii): $\so{n,\C}$ under a common action; copy~$1$ decoupled under the spectator action & fixed (restated: Prop.~13.14, Thm~17.20) & d7-1a, d7-1d\\
\cref{d7:prop:exist}(ii): $\so n$ under the interior action; $C^{12}=M_{12}+C_{\rm surp}$ & fixed & d7-1b, d7-1c\\
\cref{d7:prop:rotor}: rotor generators $\uu n=\so n\oplus\operatorname{Sym}(n)\otimes G$ & fixed; non-central rotor constructed & d7-2a\\
\cref{d7:thm:schr}(i), (ii), (iv), (v): Schr\"odinger form, unitarity iff compact, sum rules, removability & fixed (elementary) & d7-3a, d7-3b, d7-3d, d7-3e\\
\cref{d7:thm:schr}(iii): rates are gauge invariants & fixed (restated: \cref{d4:thm:covariant}) & d7-3c\\
\cref{d7:thm:schr}(vi): bilinear boundedness criterion; five cases at $n=3$ & fixed (corollary of \cref{d4:thm:criterion}) & d7-3f\\
\cref{d7:prop:boost}: copy boost, exceptional point, conserved positive form & fixed; pattern inherited & d7-4a, d7-4b\\
\cref{d7:thm:three}: spectrum, plaquette, Jarlskog determinant, reflections, $P$ and $T$, bipartite criterion, circulation & fixed; parallels inherited & d7-5a to d7-5f\\
\cref{d7:prop:isospin}: integral copy isospin; the doublet is one connection and two rate terms & fixed (restated: Prop.~10.9, \cref{d5:prop:energysym}) & d7-6a to d7-6c\\
\cref{d7:prop:so4unequal}: $\so4_\Delta\to$ diagonal $\so3$ at unequal rates & fixed; custodial parallel interpretive & d7-7a, d7-7b\\
\cref{d7:rem:verdict}: the deal breaker; the commutator; the arrow reversed & fixed; (c) constructed; (e) interpretive & d7-8\\
\bottomrule
\end{tabular}
\caption{Claims of \cref{d7:sec:main}, their classes, and the checks printed by \texttt{d7\_bfield\_between.py}.}
\label{d7:tab:claims}
\end{table}
